%; whizzy-master slides.tex
%; whizzy frame -pdf ./skimopen 

\usepackage{xspace} % To get the right spacings in front of : and so on
\usepackage[english]{babel}
\usepackage{subfigure}
\usepackage{pgf}
\usepackage{colortbl}
\usepackage{xcolor}
\usepackage{abbrevs}
\usepackage[color]{coqdoc}
\def\coqlibrary#1#2{}

\usepackage{coq}
\usepackage{tikz}
\usetikzlibrary{shapes}
\usepackage{pifont}
\usepackage{utf}
\usepackage{natbib}
\usepackage[OT1]{fontenc}
\usepackage{alltt}
\def\thetitle{\Coq 8.5: what's new, what's next?}
\def\names{M. Dénès, M. Sozeau}
\def\theauthor{\name{M. Dénès}$^{\dagger}$ \& \name{M. Sozeau}$^{*}$}
\usepackage{hyperref}


\def\vec#1{\overrightarrow{#1}}
\def\thetextitleabbr{\thetitle}
\def\thetextitle{\thetextitleabbr}
\def\thesubject{Theoretical Computer Science}
\def\thekeywords{Coq}
\def\theevent{\Coq POPL Meeting\\
January 21st 2014\\
San Diego, USA}
\def\theeventabbr{}
\def\theauthorabbr{\name{M. Dénès} \& \name{M. Sozeau}}
\usepackage{prelude-beamer}

\author{\theauthorabbr}
\def\PI{\name{PI}}
\def\shaded#1{{\color{black!50}#1}}
 
\def\seq{\fCenter}

\pgfdeclareimage[height=1.4cm]{inria-logo}{inria-logo}

\pgfdeclareimage[height=1.4cm]{coq-biglogo}{barron-logo-pdf}
\pgfdeclareimage[height=0.4cm]{coq-smalllogo}{barron-logo}
\pgfdeclareimage[height=1cm,interpolate=true]{coq-logo}{barron-logo-pdf}

\pgfdeclareimage[height=7cm]{coq-share}{coq-share}

\def\imgheight{7.5cm}

\def\typea{\vdash}
\def\suba{\rightslice}
\def\indent{\coqdocindent}
\def\var{\coqdocvar}

\def\coqdocid#1{\coqdocvar{#1}}

%\renewcommand{\Type}[1]{\kw{Type}_{#1}}

\newcommand{\elt}[4]{\ensuremath{\constr{exist}_{#1,#2}~#3~#4}}
\newcommand{\eltlight}[2]{\ensuremath{\constr{exist}~#1~#2}}
\def\class{\kw{class}~}
\def\instance{\kw{instance}~}
\def\tclass#1{\coqdocind{#1}}
\def\where{\kw{where}~}
\def\module#1{\texttt{#1}}

\def\vector#1{\ind{vector}~#1}
\def\vnil{\constr{vnil}}
\def\vcons{\constr{vcons}}

\renewcommand{\bar}[1]{{\overline{#1}}}

\setlength{\coqdocbaseindent}{0em}

\def\coqlibrary#1#2#3{}

% Tutorial

\newboolean{darkfond}
\setboolean{darkfond}{false}
%\setboolean{darkfond}{true}

\definecolor{metavarcolor}{rgb}{0.5,0.0,1.0}
\definecolor{darkgreen}{rgb}{0.1,0.7,0.1}
\definecolor{answercolor}{rgb}{1.0,0.0,0.0}
\definecolor{normalcolor}{rgb}{0.0,0.0,0.0}
\definecolor{exbluecolor}{rgb}{0.1,0.1,0.9}
\definecolor{dontlookcolor}{rgb}{0.5,0.5,0.5}
\definecolor{termcolor}{rgb}{0.0,0.1,0.9}
\definecolor{lookcolor}{rgb}{0.8,0.2,0.0}
\definecolor{prooftermcolor}{rgb}{0.3,0.1,1.0}
\definecolor{typecolor}{rgb}{1.0,0.6,0.0}
\definecolor{taccolor}{rgb}{0.70,0.10,0.0}
\definecolor{pink}{rgb}{0.8,0.6,0.6}
\definecolor{darkmagenta}{rgb}{0.4,0.0,0.6}
\definecolor{darkblue}{rgb}{0.0,0.0,0.6}

%\newenvironment{subsecframe}[2][]{\begin{frame}[environment=subsecframe,#1]\frametitle{#2}\begin{subsection}{#2}}{\end{subsection}\end{frame}}

\newenvironment{subsecframe}[1]{\begin{subsection}{#1}\begin{frame}\frametitle{#1}}{\end{frame}\end{subsection}}

\newenvironment{subsecframefragile}[1]{\begin{frame}[fragile,environment=subsecframefragile]\frametitle{#1}\begin{subsection}{#1}}{\end{subsection}\end{frame}}

\begin{document}


\setbeamertemplate{background canvas}[vertical shading][top=rouge1,middle=rouge1,bottom=rouge1]
\setbeamertemplate{footline}{\hspace{5em} \textcolor{white} {\null \hfill}\hspace{2em}\null \vspace*{3pt}}

\begin{frame}

\begin{textblock*}{40mm}[0,0](10mm,0mm)
 \includegraphics[width=6cm]{INRIA_CHERCHEURS_UK_RVB}
  \end{textblock*}

\begin{textblock*}{12cm}(13mm,50mm)
{\textcolor{white} {
{\huge \thetitle}\\[2mm]
{\theauthor}}}
\end{textblock*}


   \begin{textblock*}{40mm}[0,0](10mm,76mm)
  \begin{picture}(5,80)
\put(0,23){\includegraphics[width=4cm,height=1.5cm]{logobasrougeV1}}
\put(20,50){\tiny \textcolor{rouge2}{${}^{\dagger}$University of Pennsylvania}}
\put(20,40){\tiny \textcolor{rouge2}{${}^{*}$$πr^2$ - Inria \& PPS}}
\end{picture}
\end{textblock*}


\begin{textblock*}{7cm}(55mm,76mm)
{\textcolor{white}{{\theevent}}}
\end{textblock*}

\vspace*{-4pt}
\end{frame}


%%%%%%%%%%%%%%%%%%%%%%%%%%%

\setbeamertemplate{background canvas}[vertical shading][top=rouge1,bottom=rouge1,middle=rouge1]
\setbeamercolor{toto}{fg=blanc,bg=rouge1}

\setbeamertemplate{footline}
{
\begin{beamercolorbox}[wd=1\paperwidth,ht=15.5pt]{toto}
%\raisebox{1ex}{\includegraphics[width=25mm]{logobastrans}}
  \hspace{4em}
  \raisebox{2.5ex}
  {\thetitle{}  -- \theauthorabbr{}}\hfill 
  \raisebox{2.5ex}
  {\insertframenumber \hspace{2mm} \null }
\end{beamercolorbox}}

\setbeamertemplate{background canvas}[vertical shading][top=white,middle=white,bottom=white]

\def\bulletfail{\alert{\ding{54}}}
\def\bulletcheck{\ding{52}}

%\setbeamercolor{background}{fg=red,bg=white,text=blue}

%\setbeamercolor{subsection in toc shaded}{fg=gray,bg=white,text=blue}
\setbeamercolor{section in toc}{fg=black,bg=white}
\setbeamercolor{subsection in toc}{fg=black,bg=white}
%\setbeamertemplate{table of contents shaded}[default]
%\setbeamertemplate{subsection in toc shaded}{\textcolor{gray}}

\setbeamertemplate{section in toc shaded}
{\begin{colormixin}{20!white}{\usebeamertemplate{section in toc}}\end{colormixin}\unskip}

\setbeamertemplate{subsection in toc shaded}
{\begin{colormixin}{20!white}{\usebeamertemplate{subsection in toc}}\end{colormixin}\unskip}

\setbeamertemplate{subsubsection in toc shaded}
{\begin{colormixin}{20!white}{\usebeamertemplate{subsubsection in toc}}\end{colormixin}\unskip}

\frame<beamer>{\frametitle{\Coq 8.5}
\tableofcontents[sectionstyle=show/show,hideallsubsections]}

\section{Up and coming}
  
\begin{subsecframe}{Incremental development -- \name{E. Tassi}}
  \begin{itemize}
  \item \texttt{git} model
  \item Granularity: opaque lemmas
  \item \texttt{.vi} interfaces: cross-file compilation
  \item CoqIDE only (emacs interface hard to implement)
  \end{itemize}

  \begin{itemize}
  \item[-/++] \emph{Not optional, backwards-compatible}
  \item[+++] \emph{Faster interaction and compilation}
  \end{itemize}
\end{subsecframe}

\begin{subsecframe}{New proof engine -- \name{A. Spiwack}}
  % * new internal proof engine (perfs)
  % - eauto backtrack with ; true proof search
  % "all:", tactical "+", "once", a native "refine"

  \begin{itemize}
  \item Existential variables + \texttt{LogicT}
  \item Proof-search semantics (\texttt{;}, \texttt{+}, $\ldots$)
  \item Cleaner code: abstract monadic interface
  \item Native \texttt{refine}, dependent goals, goal manipulation
  \end{itemize}
  
  \begin{itemize}
  \item[--/++] \emph{Not optional, backwards-compatibility layer}
  \item[--] \emph{0-15\% time overhead (still invesigating, \textsc{OCaml} issue likely)}
  \end{itemize}
\end{subsecframe}

\begin{subsecframe}{Universe polymorphism -- \name{M. Sozeau}}
% Def interactive
  \begin{itemize}
  \item \emph{Checked} universes.
  \item Global definitions (and inductives) or local with prenex, ``bounded''
    quantification. \emph{Conservativity proof}
  \item \emph{Elaboration} from typical ambiguity + implicit
    generalization.
  \item Checks \texttt{HoTT/Coq}, B. Barras and T. Coquand's semi-simplicial
    model and another (1-)groupoid model.
    \pause
  \item Better error messaging, unification sensitive to universes.
  \item Generalized rewriting in \alert{\texttt{Type}}
  \end{itemize}

  \begin{itemize}
  \item[=/++] \emph{Optional except for checked
      universes, impacts the ML hacker only}. Backwards-compatibility layer.
  \item[=+/+] \emph{Comparable or better performance, more expressive}
  \end{itemize}
\end{subsecframe}

\begin{subsecframe}{Native compilation -- \name{M.~Dénès \& B.~Grégoire}}
\begin{itemize}
\item Useful for large reflection proofs
\item \Coq $"->"^{\texttt{nativecomp}}$ \textsc{ML} $"->"^{\texttt{ocamlopt}}$ \textsc{ASM}
\item Faster than \texttt{vm\_compute}
\end{itemize}

  \begin{itemize}
  \item[+] \emph{Optional}
  \item[++/--] \emph{Faster at runtime, compilation is slow}
  \end{itemize}
\end{subsecframe}

\begin{subsecframefragile}{Fast record projections -- \name{M. Sozeau}}
  \begin{itemize}
  \item Economic representation
  \item η-rule (for non-empty records only): 
    \verb|Build_rec (p1 r) .. (pn r)| is canonical
  \item Fast reduction (address of \verb|r| + offset)
  \item No more stupid unfoldings
  \end{itemize}
  
  \begin{itemize}
  \item[+/+--] \emph{Optional, backwards-compatibility layer,
      small source-level incompatibilities}.
  \item[${+}^{ω}$] \emph{Exponentially better performance}
  \end{itemize}

  \TODO{MS: Graph}

  % * Graphe de perfs
\end{subsecframefragile}

\begin{subsecframe}{New coqdoc -- \name{Y. Régis-Gianas}}
  \begin{itemize}
  \item Based on the AST and \alert{\emph{not}} lexing anymore
  \item All semantic information available
  \item Plugable interface
  \end{itemize}

  \begin{itemize}
  \item[+] \emph{Not optional but legacy \texttt{coqdoc} code included}
  \item[=] \emph{Same performance}
  \end{itemize}
  
  % \vfill
  % \begin{center}\alert{{\Huge Demo}}\end{center}
\end{subsecframe}

\begin{subsecframefragile}{Misc -- \Coq dev team \& contributors}
  \begin{itemize}
  \item Interfaces, documentation, and \textsc{OCaml} best practices
    (\name{P.M. Pédrot}, $\ldots$). \alert{Down to 8.3 performance!}
  \item Tactics in terms: \verb|$(tac)$| (\name{P.M. Pédrot})
  \item Module system simplifications (\name{P. Letouzey})
  \item Generalized intro patterns (\name{H. Herbelin, P. Letouzey})
  \item Better(?) guard condition (\name{P. Boutillier, H. Herbelin})
  \item Rewriting with strategies (\name{M. Sozeau})
  \end{itemize}

\end{subsecframefragile}

\begin{subsecframe}{\texttt{opam} -- \name{T. Braibant}}
  
  \begin{itemize}
  \item \name{Coq, HoTT/Coq, Coq+Ssreflect}, any version, git included
  \item Some user contribs: \name{Containers, Coccinelle, Ergo, \dots}
  \item Submit a pull request! Try! Test!
  \end{itemize}

  \begin{center}
    \url{http://opam.ocaml.org/}
    \vspace{0.5em}

    \url{https://github.com/braibant/opam-coq-repo}
  \end{center}
  \TODO{TB: USB key?}
\end{subsecframe}

\begin{subsecframe}{User libraries, plugins}

  \begin{itemize}
  \item \name{Ssreflect} 1.5 (compatibility with TCs, split between
    \name{ssr} and \name{math-comp}, more theories: \url{http://www.msr-inria.fr/projects/mathematical-components/})
  \item \name{Mtac}, \name{Cybele}: ``internalized'' tactic language
  \item \name{Paco}: coinduction, logically
  \item \name{Extlib}, \name{TLC}: extensions of the standard library
  \item \name{Why3}: call SMT solvers from \Coq
  \end{itemize}
  
\end{subsecframe}

\begin{subsecframe}{The bug}
  \begin{itemize}
  \item Due to a (very) extensional axiom, related to \emph{equality} on
    types interacting with the guard condition.
  \item A good fix is known, but breaks the stdlib's use of the smart
    guard condition (ha ha). All definitions can be adapted. Work by
    \name{B. Barras, C. Paulin, M. Dénès, P. Boutilier}, $\ldots$
  \item[=] Happens in \name{Agda} too!
  \item Argues for less syntax and more logic.
  \end{itemize}
\end{subsecframe}

\begin{subsecframe}{Tentative planning}
  \begin{itemize}
  \item February 2014: 8.5-α release
  \item April 2014: 8.5-β release and opening of new branch
    (only bug fixes in trunk between α and β)
  \item August 2014: final release
  \end{itemize}
\end{subsecframe}


\section{The close future}
\frame<beamer>{\frametitle{The close future}
\tableofcontents[sectionstyle=show/show,hideallsubsections]}

\begin{subsecframe}{Performance}

  \begin{itemize}
  \item Caching of typechecked terms (\name{E. Tassi, M. Sozeau})
  \item Generalized hashconsing (\name{T. Braibant, M. Sozeau,
      P. M. Pédrot, $\ldots$})
  \end{itemize}
  
\end{subsecframe}

\begin{subsecframe}{CoqMT -- \name{P.Y. Strub} et al}
  \begin{itemize}
  \item Extend the definitional equality with a (first-order) decidable
    theory, e.g. Presburger arithmetic, constructors
  \item[$"->"$] A \emph{strict} equality avoiding coercions, a slightly
    larger TCB (SN proof available).
  \item Some limitations on the use of equalities in context
  \item Current state: useful but not as well-integrated as it could
    be, still at 8.3.
  \item To be integrated as an optional feature
  \end{itemize}
\end{subsecframe}  

\section{The distant future}
\frame<beamer>{\frametitle{The distant future}
\tableofcontents[sectionstyle=show/show,hideallsubsections]}

\begin{subsecframe}{HoTT or OTT?}
  HoTT ideas, math oriented:
  \begin{itemize}
  \item HITs (prototype by \name{B. Barras})
  \item Resizing rules
  \item \alert{\emph{Definitionally}}: βδ, η?, ι? 
  \item \alert{\emph{Propositionally}}: proof-irrelevance, functional extensionality,
    isomorphic types are equal (from an axiom).
  \end{itemize}
  
  OTT/Epigram ideas, CS oriented:
  \begin{itemize}
  \item Levitation of inductive types
  \item Clear $\Prop$/$\Set$ separation
  \item \alert{\emph{Definitionally}}: βδηι, proof-irrelevance, functional
    extensionality.
  \item \alert{\emph{Propositionally}}: defined, heterogeneous on
    values, structural on types ($\neq$ iso). 
    With coercion and coherence operators.
  \end{itemize}

  Ideally, we should mix the two\dots
\end{subsecframe}

\begin{subsecframe}{\Coq as a programming language}

  Usability, effiency woes:
  \begin{itemize}
  \item I/O, FFI (\TODO{MD: Demo Hello world! in Coq})
  \item A sliding scale for efficiency/safety in the kernel (ideas:
    tainting code/data, signatures). Allowing, unchecked recursive
    definitions, custom reductions\dots
  \item Definitional proof-irrelevance based on bracket types (useful
    for DTP and recursion encodings)
  \item Type-based termination (sized-types, \name{J. L. Sachini})
  \end{itemize}
\end{subsecframe}

% \begin{subsecframe}{Consortium}
% \end{subsecframe}
% \begin{subsecframe}{Development model}
%   Recentralize (\name{P. Letouzey}, \Coq dev team)
%   \begin{itemize}
%   \item , bugzilla, website on new machine (TODO by Pierre L), patches, dependencies of Coq, publishing libraries TODO : move to discussion ?

\end{document}

%%% Local Variables: 
%%% mode: latex
%%% TeX-PDF-mode: t
%%% TeX-master: "slides"
%%% End: 
